% Modern constructions filling presentation/adequacy components, not original-source replacements.
\section{Exact minimal and intermediate presentations}
\label{sec:p08-minimal-and-intermediate-presentations}

Minimal logic, KC and LC require different semantic restrictions. Minimal
logic removes falsity elimination; its falsum may hold at a world without
every formula holding there. KC keeps intuitionistic falsum and adds weak
excluded middle; LC keeps intuitionistic falsum and adds prelinearity.
This section supplies exact modern propositional calculi and complete
local proofs. It does not identify these constructions with Johansson's
primitive-negation presentation, Jankov's original pseudo-Boolean-algebra
argument, or Dummett's particular denumerable matrix
\cite[Section~3.1]{ColacitoDeJonghVargas2016SubminimalDraft}
\cite{Jankov1968WeakExcludedMiddle,Dummett1959}.
The wider intermediate lattice remains
Definition~\ref{def:p08-intermediate-lattice} and the separate proved
algebraic correspondence in Theorem~\ref{thm:p10-intermediate-variety-correspondence}.

\begin{hbdefinition}{def:p08-mi-language}{The exact four propositional derivability relations}
Fix any set of atoms. Formulas are finite expressions over those atoms and
\(\bot,\land,\lor,\to\); \(\neg A\) abbreviates \(A\to\bot\).
The notation \(\Gamma\derives_{\mathsf{MPL}}A\) means a finite tree on
Sheet~\ref{sheet:minimal-nd} with conclusion \(A\) and every actual open
assumption occurrence in \(\Gamma\). Every selected discharge class may
be empty. There are no quantifiers, equality rules, RAA or falsity
elimination. This is a formal presentation, not an abbreviation for every
system called minimal logic.

Write \(\mathsf H_0\) for the common nine-schema intuitionistic base on
Sheets~\ref{sheet:hilbert-kc} and~\ref{sheet:hilbert-lc}, excluding each
sheet's additional named schema and retaining modus ponens. The judgments
\(\Gamma\derives_{\mathsf H_0}A\),
\(\Gamma\derives_{\mathsf H_{\mathsf{KC}}}A\) and
\(\Gamma\derives_{\mathsf H_{\mathsf{LC}}}A\) mean finite formula
sequences of premises, displayed schema instances and modus-ponens
conclusions. A schema instance has arbitrary formulas substituted for its
metavariables. There is no rule substituting inside a fixed premise while
leaving its premise context unchanged. In particular all four relations
have finite premise support even when \(\Gamma\) is infinite.
\index{minimal logic!defined-negation propositional presentation}
\index{KC!exact Hilbert presentation}\index{LC!exact Hilbert presentation}
\end{hbdefinition}

\begin{hbdefinition}{def:p08-mi-models}{Minimal falsum and the two intermediate frame classes}
A minimal model consists of a nonempty partial order \((W,\leq)\), a
persistent atom valuation \(V\), and an upward-closed set \(F_{\bot}\subseteq W\).
Forcing an atom means membership in its valuation; forcing \(\bot\) means
membership in \(F_{\bot}\). Conjunction and disjunction have their
same-world clauses, and \(w\forces A\to B\) means that every \(v\geq w\)
forcing \(A\) forces \(B\). Thus \(w\forces\neg A\) means that each
future \(A\)-world lies in \(F_{\bot}\), not that no future world forces
\(A\). This is the falsum formulation compared with the selected modern
source; its independent local proof is below.

For KC and LC set \(F_{\bot}=\varnothing\), exactly as in the IPL model
definition. A KC frame here is a rooted partial order with a greatest world;
it need not be finite. An LC frame is a rooted linearly ordered partial
order; it need not have a greatest world. Consequence for each of these
three classes is local: every model and its root forcing all members of
\(\Gamma\) must force the conclusion. Restricting a model to a future cone
does not change truth at its worlds. The common-world premise regime is
part of each definition.
\end{hbdefinition}

\begin{hblemma}{thm:p08-mi-derivation-operations}{proved=annex:p8.mi.operations}{Deduction, premise replacement, and the intuitionistic base bridge}
Each relation in Definition~\ref{def:p08-mi-language} is monotone in
premises, has finite character, admits replacement of finitely many derived
premises by their derivations, and satisfies
\[
 \Gamma\cup\{A\}\derives B\quad\Longleftrightarrow\quad
 \Gamma\derives A\to B
\]
in that same calculus. Each admits conjunction introduction/projections,
disjunction introductions and reasoning by cases. The Hilbert relations
also admit explosion; MPL does not import it.
For every \(\Gamma,A\),
\[
 \Gamma\derives_{\mathsf H_0}A
 \quad\Longleftrightarrow\quad \Gamma\derives_{\mathsf{IPL}}A,
\]
where IPL means exactly the propositional LJ relation of
Definition~\ref{def:p08-ipl-completeness-setting}. Adjoining falsity
elimination to the MPL sheet gives the propositional empty-discharge NJ
presentation, whose consequence relation is the same IPL relation.
\end{hblemma}

\begin{hbtheorem}{thm:p08-mi-minimal-adequacy}{proved=annex:p8.mi.minimal}{Soundness and strong completeness of the selected minimal calculus}
In classical ZFC, for any atom set, premise set \(\Gamma\), and formula
\(A\), the MPL sheet derives \(A\) from \(\Gamma\) exactly when every
minimal model rooted at a world forcing \(\Gamma\) forces \(A\).
When derivability fails, a set-sized rooted partial-order countermodel
exists, with a persistent falsum upset that may be nonempty.
The result is propositional; it adds no first-order or equality theorem.
\end{hbtheorem}

\begin{hbexample}{def:p08-mi-abnormal-world}{Falsum does not license explosion in minimal logic}
Take one world \(r\), let \(F_{\bot}=\{r\}\), and let \(V(p)=\varnothing\).
Then \(r\forces\bot\), \(r\not\forces p\), and
\(r\not\forces\bot\to p\). For every formula \(B\), however,
\(r\forces B\to\bot\), since its only future world forces falsum.
Thus forcing a negation at such a world is not evidence that its operand
is absent. Soundness proves both \(\{\bot\}\not\derives_{\mathsf{MPL}}p\)
and the non-theoremhood of explosion. Setting falsum false everywhere
would exclude this required separating model.
\end{hbexample}

\begin{hbtheorem}{thm:p08-mi-minimal-falsum-recovery}{proved=annex:p8.mi.falsum-recovery}{Falsum recovery and the positive-language bridge}
For every formula \(B\), MPL derives each direction between \(\bot\)
and \(\neg B\land\neg\neg B\). In particular its falsum worlds are
exactly those worlds forcing both some formula and its defined negation.
Choose a new distinguished atom \(f\) outside the original atom set,
and replace \(\bot\) by \(f\), commuting with the three binary
connectives. MPL derivability is equivalent, under that translation of
premises and conclusion, to the positive Hilbert base obtained from
\(\mathsf H_0\) by omitting the Falsity schema and taking its language
to be atoms and \(\land,\lor,\to\). This is the defined-negation
falsum/positive-language bridge; it is not a translation theorem for
an independently primitive-negation language.
\end{hbtheorem}

\begin{hbtheorem}{thm:p08-mi-intermediate-adequacy}{proved=annex:p8.mi.intermediate}{Strong frame completeness for the exact KC and LC sheets}
In classical ZFC, for any atom set, any \(\Gamma\), and any \(A\),
\(\Gamma\derives_{\mathsf H_{\mathsf{KC}}}A\) iff \(A\) follows locally
from \(\Gamma\) on the rooted greatest-world models of
Definition~\ref{def:p08-mi-models};
\(\Gamma\derives_{\mathsf H_{\mathsf{LC}}}A\) iff it follows locally
on that definition's rooted linearly ordered models. In either failure
case there is a set-sized countermodel in the named class.
These are modern frame proofs for these exact sheets, not a promotion
of the separately cited original-source profiles to source-native proofs.
\end{hbtheorem}

\begin{hbtheorem}{thm:p08-mi-finite-decision}{proved=annex:p8.mi.finite}{Finite countermodel bounds and terminating decision}
For any of MPL, the exact KC sheet or the exact LC sheet, suppose finite
\(\Gamma\) does not derive \(A\). Expand defined negations and let
\(\Sigma\) be the set of subformulas in \(\Gamma\cup\{A\}\).
There is a rooted countermodel of the named kind with at most
\(2^{|\Sigma|}\) worlds. Consequently finite-input consequence and
theoremhood are decidable, assuming finite effective formula encoding and
decidable equality of atom names. No optimal complexity bound is asserted.
For minimal models the search includes all persistent falsum upsets;
for KC and LC it fixes that upset empty and retains the specified frame
restriction.
\end{hbtheorem}

\begin{hbmethod}{def:p08-mi-decision-method}{An explicit finite search}
For sizes from one through the preceding bound, enumerate every relation
on that finite labeled set and keep just the rooted partial orders.
For KC additionally require a greatest element; for LC require linearity.
Enumerate every upward-closed valuation of the finitely many input atoms,
and, for MPL, every upward-closed falsum set. Evaluate formulas by
structural recursion, testing every future world in an implication clause.
A root satisfying the premises and failing the conclusion is a negative
answer with a concrete countermodel. If the finite search finds none, the
finite-countermodel theorem gives a positive answer. Proof enumeration
over these finite effective calculi can also recover a derivation for
that positive answer; an arbitrary bounded proof search is not the
decision argument.
\end{hbmethod}

\begin{hbtheorem}{thm:p08-mi-separations}{proved=annex:p8.mi.separations}{Exact separations, disjunction failures, and the finite greatest-world condition}
With the shared formula grammar and theorem-set reading,
\[
 \mathsf{MPL}\subsetneq\mathsf{IPL}\subsetneq
 \mathsf{KC}\subsetneq\mathsf{LC}\subsetneq\mathsf{CPC}.
\]
Here the KC and LC theorem sets are those of the exact new sheets.
KC and LC both fail the theorem-level disjunction property.
For a finite rooted partial order, weak excluded middle is valid under
every persistent valuation iff the order has a greatest world.
Having a unique maximal world is equivalent to that condition only under
the finite hypothesis; the equivalence fails for infinite rooted orders.
\end{hbtheorem}

\begin{hbdefinition}{def:p08-mi-new-bottom-algebra}{The finite new-bottom algebra}
Let \(\alpha\) be a finite Heyting algebra, with its own bottom and
top. Form \(\alpha^*\) by adjoining a new element \(f\) strictly
below every old element, retaining the old lattice order and extending
meet and join by \(f\wedge a=f\), \(f\vee a=a\).
Its bottom is the new \(f\), its top is the old top, and its conditional
is the relative pseudocomplement in this finite distributive lattice.
The letters \(\alpha,\alpha^*,f\) here denote algebras and an algebra
element, not the distinguished syntactic atom \(f\) of the previous
translation. This is the new-bottom construction used in Jankov's
opening algebraic presentation \cite[pp.~997--999]{Jankov1968WeakExcludedMiddle};
the following independent bridge derives its adequacy from the local
frame result rather than replacing his proof or its further variants.
\end{hbdefinition}

\begin{hbtheorem}{thm:p08-mi-kc-algebra-bridge}{proved=annex:p8.mi.kc-algebras}{KC and the finite new-bottom algebras}
For finite \(\Gamma\) and \(A\), the KC sheet derives \(A\) from
\(\Gamma\) iff every valuation in every \(\alpha^*\) of
Definition~\ref{def:p08-mi-new-bottom-algebra} assigns \(A\) top
whenever it assigns every premise top. This includes the one-element
old algebra, whose new-bottom extension is the two-element algebra.
Thus the theoremhood statement is a derived algebraic bridge to the
selected source construction, with premise preservation proved only
for the finite-input scope stated here.
\end{hbtheorem}

\begin{hbtheorem}{thm:p08-mi-finite-chain-bridge}{proved=annex:p8.mi.chains}{Finite-input LC and numerical Godel chains}
For finite \(\Gamma\) and \(A\), LC derives \(A\) from \(\Gamma\)
iff every valuation in every finite bounded chain assigns \(A\) its top
value whenever it assigns every premise its top value, using minimum,
maximum, bottom, and the conditional operation
\[
 a\Rightarrow b=\begin{cases}1,&a\leq b,\\ b,&a>b.\end{cases}
\]
The same finite-input equivalence holds for any fixed infinite bounded
linearly ordered set with these operations. In this theorem
\(0,1,\Rightarrow\) are algebraic bounds and an algebraic operation,
not a falsum-world set or a sequent arrow. No arbitrary-premise result
for a fixed numerical chain, nor identification with Dummett's particular
denumerable matrix or Godel's 1932 finite matrices, is asserted.
\end{hbtheorem}

The frame/algebra bridge just proved is independent of the source-specific
matrix choices. Jankov's original algebraic completeness variants and
Dummett's selected denumerable-matrix derivation remain separate
presentations \cite{Jankov1968WeakExcludedMiddle,Dummett1959}.
Similarly the source's primitive-negation minimal language and its
translation to a falsum presentation require their own stated translation;
changing a connective grammar is not literal proof identity.
